% ---------------------------------------------------------------------------
% Zusammenfassung / Fazit (Merkblatt 3.2: max. ~15% der Seiten, keine neuen
% Argumente; abschliessend ggf. Fazit / Ausblick)
% ---------------------------------------------------------------------------
\chapter{Zusammenfassung und Fazit}
\label{cha:fazit}


\section{Zusammenfassung}
\label{sec:zusammenfassung}


Diese Arbeit reproduzierte die Studie von \textcite{Kanas:2023} mit dem Titel „Greater local supply of language courses improves refugees’ labor market integration“ und analysierte die Variable „Asylstatus“ als Interaktionseffekt. Diese Arbeit versuchte daher, die folgende Forschungsfrage zu beantworten: Profitieren Geflüchtete mit anerkanntem Asylantrag stärker von lokalen Sprachkursen im Hinblick auf ihren Erfolg auf dem Arbeitsmarkt?

 \textcite{Kanas:2023} haben den kausalen Einfluss des Angebots an lokalen Integrationskursen in Deutschland auf die Integration von Geflüchteten in den Arbeitsmarkt analysiert. Diese Studie stützt sich auf den Datensatz zur Geflüchtetenwelle, die zwischen 2016 und 2019 nach Deutschland kam. Diese Arbeit hat die Analyse von \textcite{Kanas:2023} unter Verwendung der Datensätze „IAB-BAMF-SOEP-Befragung von Geflüchteten“, „Kreisebene-Makrodaten“ und „SOEP Regional Data“ reproduziert. Zusätzlich zur Reproduktion wurde die Analyse um den Asylstatus als Moderator erweitert.
 
 Im Folgenden werden einige zentrale Ergebnisse der Reproduktionsstudie vorgestellt. Zunächst liegen Ergebnisse zum Einfluss des Angebots an lokalen Integrationskursen auf die Beschäftigung vor. Ein Anstieg des Angebots an lokalen Integrationskursen um eine Standardabweichung erhöht die Beschäftigungswahrscheinlichkeit von Geflüchteten um etwa 2,98 Prozentpunkte. Die Wechselwirkung zwischen dem lokalen Kursangebot und der Aufenthaltsdauer führt zu einem positiven Ergebnis, ist jedoch statistisch nicht signifikant. Daher gibt es keine Belege dafür, dass sich die Auswirkungen des lokalen Integrationskursangebots auf die Beschäftigung in Abhängigkeit von der Aufenthaltsdauer ändern.
 
 Das Ergebnis der Reproduktion zeigte auch, dass das Kursangebot mit bestimmten intermediären Ergebnissen in Zusammenhang steht. Das lokale Kursangebot wies einen positiven Zusammenhang mit den Deutschkenntnissen, dem Abschluss eines Integrationskurses und dem Sprachnachweis auf. Es wurde jedoch ermittelt, dass dieser Effekt mit zunehmender Aufenthaltsdauer abnimmt. Schließlich wurde festgestellt, dass das lokale Kursangebot keinen signifikanten Einfluss auf die Häufigkeit der Kontakte zu Deutschen hatte. Damit stimmen die Ergebnisse der Reproduktion mit den zentralen Befunden von \textcite{Kanas:2023} überein.
 
 Im Rahmen dieser Erweiterung wurde analysiert, ob sich die Auswirkungen des Kursangebots auf die Beschäftigung je nach Asylstatus der Geflüchteten unterscheiden. Die Variable „Aufenthaltsstatus“ wird, wie von \textcite{Kanas:2023} definiert, in drei Gruppen unterteilt: „keine Entscheidung“, „anerkannt“ und „abgelehnt“. Die Forschungsfrage wurde anhand des Spracherwerbsmodells von \textcite{Esser:2006} erläutert. Nach \textcite{Esser:2006} basiert der Spracherwerb auf rationalen Entscheidungen. Aus der Perspektive der Bleibeperspektive wurde angenommen, dass Geflüchtete mit einem sichereren Aufenthaltsstatus tendenziell mehr in den Spracherwerb investieren. Eine weitere Annahme basierte auf dem ungewissen Status der Geflüchteten und den rechtlichen Hindernissen beim Zugang zu Sprachkursen. Da Geflüchtete mit ungewissem Status keinen uneingeschränkten Zugang zu Sprachkursen haben, wurde angenommen, dass sie stärker von Sprachkursen auf Kreisebene profitieren könnten.
 
 Der Zusammenhang zwischen dem Kursangebot und der Beschäftigung zeigt keine Unterschiede je nach Asylstatus. Daher wird die Annahme, dass anerkannte Geflüchtete stärker vom lokalen Kursangebot profitieren, nicht gestützt. Als Erweiterung wurden zwei verschiedene Messungen des Asylstatus durchgeführt. Die erste betraf den Asylstatus im Erhebungsjahr (E1.1), die zweite den Asylstatus bei der ersten Beobachtung (E1.2). Die Erweiterungsanalyse hat gezeigt, dass der Zusammenhang zwischen dem lokalen Kursangebot und der Beschäftigung je nach Asylstatus keinen statistisch signifikanten Unterschied aufweist. Dieser Befund lässt sich bei beiden Arten der Messung des Asylstatus beobachten (E1.1 und E1.2).
 
 Die Erweiterung für intermediäre Mechanismen wurde zusätzlich getestet. Als intermediäre Mechanismen wurden die deutschen Sprachkenntnisse, der Besuch eines Integrationskurses, der Abschluss eines Integrationskurses, ein Sprachzertifikat sowie Kontakte zu Deutschen untersucht. Es konnte jedoch bei keinem dieser Mechanismen ein statistisch signifikanter Unterschied hinsichtlich des Asylstatus der Geflüchteten festgestellt werden. In der explorativen Analyse wurde geprüft, ob sich die Auswirkungen des Angebots an Integrationskursen auf die Erwerbsbeteiligung je nach Aufenthaltsdauer und Asylstatus unterscheiden. Es konnten jedoch keine Belege dafür gefunden werden, dass sich die Auswirkungen des Kursangebots auf die Erwerbsbeteiligung je nach Aufenthaltsdauer und Asylstatus unterscheiden.
 
 
 
 \section{Fazit}
 \label{sec: fazit}
 
 Das Ziel dieser Arbeit ist es, den Zusammenhang zwischen lokalen Integrationskursen und der Beschäftigung von Geflüchteten zu reproduzieren und als Beitrag zu zeigen, ob sich dieser Zusammenhang je nach Asylstatus unterscheidet. Die Ergebnisse der Reproduktion bestätigen die zentrale Erkenntnis von \textcite{Kanas:2023}. Die Reproduktion zeigte, dass ein positiver Zusammenhang zwischen dem lokalen Kursangebot und den Beschäftigungsaussichten von Geflüchteten besteht. Der Asylstatus liefert keine Hinweise darauf, dass er die Haupterkenntnis beeinflusst. Es gibt keine Belege dafür, dass sich der Zusammenhang zwischen dem lokalen Kursangebot und der Beschäftigung von Geflüchteten je nach Asylstatus unterscheidet. Dieses Ergebnis bleibt auch bei beiden unterschiedlichen Messungen des Asylstatus bestehen. Daher gibt es keinen Befund, der die Annahme stützt, dass anerkannte Flüchtlinge stärker von lokalen Sprachkursen profitieren würden. Berücksichtigt man Mechanismen wie Kursabschluss und Deutschkenntnisse, bleibt der Einfluss des Kursangebots auf die Beschäftigung weiterhin bestehen. Es lässt sich jedoch keine Differenzierung nach dem Asylstatus beobachten.
 Die Erwartung, dass anerkannte Geflüchtete stärker von lokalen Sprachkursen profitieren, wurde mit dem Spracherwerbsmodell von \textcite{Esser:2006} begründet. In den Erweiterungsanalysen wurde jedoch kein empirischer Befund ermittelt, der diese Annahme bestätigt. In den Erweiterungsanalysen zeigte sich jedoch kein statistisch signifikanter Unterschied in der Wirkung des lokalen Kursangebots auf die Erwerbstätigkeit zwischen den Asylstatusgruppen. Wenn die Stichprobe jedoch in drei Gruppen nach Asylstatus unterteilt wird, verringern sich die Gruppengrößen. Bei kleinen Gruppengrößen kann es schwieriger sein, tatsächlich bestehende kleinere oder moderate Unterschiede statistisch nachzuweisen. Dass kein signifikanter Unterschied besteht, bedeutet daher nicht zwangsläufig, dass ein theoretisch vorhergesagter Zusammenhang tatsächlich nicht vorliegt.
 
 Die Zuweisung zu einem Kreis erfolgt zufällig, weshalb die Hauptergebnisse von \textcite{Kanas:2023} kausal interpretiert werden können. Der Aufenthaltsstatus wird nicht zufällig festgelegt, sondern richtet sich nach bestimmten rechtlichen Vorgaben. Außerdem wird der Aufenthaltsstatus erst nach der Zuweisung zu einem Kreis festgelegt. Aus diesen Gründen lassen sich die Ergebnisse der „Erweiterung“ korrelativ interpretieren. Zukünftige Forschungsarbeiten könnten sich mit diesen Einschränkungen auseinandersetzen. Durch eine größere Stichprobe könnten möglicherweise auch kleinere Unterschiede zwischen den Asylstatus-Kategorien statistisch nachgewiesen werden.
 
 
 
 
 
 
 
 
 
 
 
 
 
 
 
 
 